\section{Auswertung}
\label{sec:Auswertung}

\subsection{Berechnung des Saugvermögens mit der Evakuierungskurve für die Drehschieberpumpe}
Um das Saugvermögen mithilfe der Evakuierungskurve zu bestimmen wird \autoref{eqn:100} 
genutzt. Es wird für jeden gemessenen Druck $p$, 
sowie den Anfangsdruck $p_0$ und den Enddruck $p_E = \qty{0.02}{\milli\bar}$ 
\begin{equation}
  \label{eqn:100}
  -\frac{S_\text{eff}}{V} \cdot t = \ln\left(\frac{p(t) - p_e}{p_0 - p_e}\right)
\end{equation}
berechnet, um daraus durch Umstellen auf das Saugvermögen der Pumpe schließen zu
können. Der Logarithmus in Abhängigkeit der Zeit ist in \autoref{fig:10} mit den
Ausgleichsgeraden dargestellt. Es wird unterteilt in drei verschiedenen Abschnitte, 
zu denen lineare Regressionen durchgeführt werden. Das Intervall (0-150), (150-290)
und (290-599). Die dazugehörigen Messwerte befinden sich in \autoref{tab:10}.
\begin{figure}
  \centering
  \includegraphics{ersteKurve.pdf}
  \caption{Darstellung des Logarithmus in Abhängigkeit der Zeit mit Ausgleichsgeraden}
  \label{fig:10}
\end{figure}
Es ergeben sich aus den Regressionen die steigungsparameter der Ausgleichsgeraden
\begin{align}
  a_1 =& \qty{-0.04107 +- 0.00015}{}\\
  a_2 =& \qty{-0.01443 +- 0.00022}{}\\
  a_3 =& \qty{-0.00587 +- 0.00002}{}
\end{align}

Bei Multiplikation mit dem Volumen $\qty{34 \pm 3.4}{\litre}$, ergeben sich die 
Saugvermögen in den jeweiligen Bereichen zu
\begin{align}
  S_\text{eff1}^\text{eva} =& \qty{4.9+-0.5}{m^3/h}\\
  S_\text{eff2}^\text{eva} =& \qty{1.71+-0.17}{m^3/h}\\
  S_\text{eff3}^\text{eva} =& \qty{0.70+-0.07}{m^3/h}
\end{align}



\begin{table}[H]
\centering
\caption{Messdaten der Evakuierungskurve der Drehschieberpumpe}
\label{tab:10}
\begin{tabular}{|r|r@{\,$\pm$\,}l|r@{\,$\pm$\,}l|p{0.5cm}|r|r@{\,$\pm$\,}l|r@{\,$\pm$\,}l|}
\hline
t / s & \multicolumn{2}{c|}{p/\unit{\milli\bar}} & \multicolumn{2}{c|}{$\log{\frac{p-p_E}{p_0-p_E}}$} & & t/s & \multicolumn{2}{c|}{p / \unit{\milli\bar}} & \multicolumn{2}{c|}{$\log{\frac{p-p_E}{p_0-p_E}}$} \\
\hline
1   & 1006   & 36    & -0.001 & 0.036 & & 301 & 0.320  & 0.032 & -8.12  & 0.11  \\
11  & 913    & 36    & -0.098 & 0.039 & & 311 & 0.300  & 0.030 & -8.19  & 0.11  \\
21  & 576    & 36    & -0.559 & 0.063 & & 321 & 0.290  & 0.029 & -8.22  & 0.11  \\
31  & 392    & 36    & -0.944 & 0.092 & & 331 & 0.260  & 0.026 & -8.34  & 0.11  \\
41  & 267    & 36    & -1.33  & 0.13  & & 341 & 0.250  & 0.025 & -8.38  & 0.11  \\
51  & 176    & 36    & -1.74  & 0.20  & & 351 & 0.230  & 0.023 & -8.48  & 0.11  \\
61  & 117    & 36    & -2.16  & 0.31  & & 361 & 0.220  & 0.022 & -8.52  & 0.11  \\
71  & 77     & 36    & -2.57  & 0.47  & & 371 & 0.210  & 0.021 & -8.58  & 0.11  \\
81  & 48     & 36    & -3.04  & 0.74  & & 381 & 0.200  & 0.020 & -8.63  & 0.11  \\
91  & 31     & 36    & -3.5   & 1.2   & & 391 & 0.190  & 0.019 & -8.69  & 0.11  \\
101 & 20     & 36    & -3.9   & 1.8   & & 401 & 0.180  & 0.018 & -8.75  & 0.11  \\
111 & 13     & 36    & -4.4   & 2.8   & & 411 & 0.170  & 0.017 & -8.81  & 0.11  \\
121 & 9.10   & 0.91  & -4.71  & 0.10  & & 421 & 0.160  & 0.016 & -8.88  & 0.11  \\
131 & 5.90   & 0.59  & -5.14  & 0.10  & & 431 & 0.150  & 0.015 & -8.95  & 0.12  \\
141 & 4.00   & 0.40  & -5.53  & 0.10  & & 441 & 0.140  & 0.014 & -9.03  & 0.12  \\
151 & 2.90   & 0.29  & -5.86  & 0.10  & & 451 & 0.140  & 0.014 & -9.03  & 0.12  \\
161 & 2.10   & 0.21  & -6.18  & 0.10  & & 461 & 0.130  & 0.013 & -9.12  & 0.12  \\
171 & 1.70   & 0.17  & -6.40  & 0.10  & & 471 & 0.120  & 0.012 & -9.22  & 0.12  \\
181 & 1.30   & 0.13  & -6.67  & 0.10  & & 481 & 0.120  & 0.012 & -9.22  & 0.12  \\
191 & 1.10   & 0.11  & -6.84  & 0.10  & & 491 & 0.110  & 0.011 & -9.32  & 0.12  \\
201 & 0.930  & 0.093 & -7.01  & 0.10  & & 501 & 0.110  & 0.011 & -9.32  & 0.12  \\
211 & 0.810  & 0.081 & -7.15  & 0.10  & & 511 & 0.100  & 0.010 & -9.44  & 0.13  \\
221 & 0.700  & 0.070 & -7.30  & 0.10  & & 521 & 0.0980 & 0.0098 & -9.47 & 0.13  \\
231 & 0.630  & 0.063 & -7.41  & 0.10  & & 531 & 0.0950 & 0.0095 & -9.50 & 0.13  \\
241 & 0.570  & 0.057 & -7.51  & 0.10  & & 541 & 0.0910 & 0.0091 & -9.56 & 0.13  \\
251 & 0.510  & 0.051 & -7.63  & 0.10  & & 551 & 0.0880 & 0.0088 & -9.60 & 0.13  \\
261 & 0.450  & 0.045 & -7.76  & 0.10  & & 561 & 0.0850 & 0.0085 & -9.65 & 0.13  \\
271 & 0.420  & 0.042 & -7.83  & 0.11  & & 571 & 0.0820 & 0.0082 & -9.70 & 0.13  \\
281 & 0.380  & 0.038 & -7.94  & 0.11  & & 581 & 0.0780 & 0.0078 & -9.76 & 0.13  \\
291 & 0.360  & 0.036 & -7.99  & 0.11  & & 591 & 0.0760 & 0.0076 & -9.80 & 0.14  \\
\hline
\end{tabular}
\end{table}

\subsection{Bestimmung des Saugvermögens der Drehschieberpumpe durch Leckratenmethode}
Die Messwerte des Druckanstiegs nach Abschiebern des Pmpenventiels sind in den
Tabellen \autoref{tab:11},\autoref{tab:12},\autoref{tab:13},\autoref{tab:14} gegeben. 
Aus diesen Drücken wird durch 
\begin{equation}
  \label{eqn:20}
  S_\text{eff}^\text{leck} = \frac{V}{p_\text{g}} \cdot \frac{\increment p}{\increment t}
\end{equation}
und mithilfe der Steigungsparameter der Ausgleichsrechnungen bei den jeweiligen
Gleichgewichtsdrücken, welche dem Quotienten $\frac{\increment p}{\increment t}$
entsprechen, das Saugvermögen bestimmt. In den folgenden Abbildungen sind die
Messwerte inklusive Ausgleichsgeraden geplottet. Es ist zu erkennen, dass einige
erste Werte bei der Regression ausgeschlossen wurden, sofern sie stark von einem
linearen verlauf abgewichen haben. Diese starken Abweichungen am Anfang liegen
wohlmöglich an eingeschlossener Luft im Ventil, welche schlagartig beim öffnen
frei wird.

\begin{figure}[H]
  \includegraphics{"drehschieberLeckrate1.pdf"}
\end{figure}
\noindent Der Steigungsparameter der Ausgleichsgeraden beträgt bei einem 
Gleichgewichtsdruck von $p_\text{g} = \qty{0.502}{\milli\bar}$
\begin{equation}
  m_{\qty{0.502}{\milli\bar}} = \qty{0.00280 +- 0.00002}{\milli\bar\per\second}.
\end{equation}

\begin{figure}[H]
  \includegraphics{"drehschieberLeckrate2.pdf"}
\end{figure}
\noindent Bei der zweiten Messung mit $p_\text{g} = \qty{10}{\milli\bar}$ ergibt 
sich
\begin{equation}
  m_{\qty{10}{\milli\bar}} = \qty{0.29483 +- 0.00040}{\milli\bar\per\second}.
\end{equation}

\begin{figure}[H]
  \includegraphics{"drehschieberLeckrate3.pdf"}
\end{figure}
\noindent Bei $p_\text{g} = \qty{50.1}{\milli\bar}$ beläuft sich der Steigungsparameter 
auf
\begin{equation}
  m_{\qty{50.1}{\milli\bar}} = \qty{2.11171 +- 0.00216}{\milli\bar\per\second}.
\end{equation}

\begin{figure}[H]
  \includegraphics{"drehschieberLeckrate4.pdf"}
\end{figure}
\noindent Zuletzt ergibt sich bei $p_\text{g} = \qty{105}{\milli\bar}$ ein Parameter von
\begin{equation}
  m_{\qty{105}{\milli\bar}} = \qty{3.86502 +- 0.00826}{\milli\bar\per\second}.
\end{equation}


\noindent Aus diesen Steigungen der Ausgleichsrechnungen wird über \autoref{eqn:20}
das Saugvermögen bestimmt als:
\begin{align}
  S_\text{eff,0.5mbar}^\text{leck} =& \qty{0.68+-0.07}{\meter\cubed\per\hour}\\
  S_\text{eff,10mbar}^\text{leck} =& \qty{3.6+-0.4}{\meter\cubed\per\hour}\\
  S_\text{eff,50mbar}^\text{leck} =& \qty{5.1+-0.5}{\meter\cubed\per\hour}\\
  S_\text{eff,105mbar}^\text{leck} =& \qty{4.5+-0.5}{\meter\cubed\per\hour}
\end{align}


\begin{table}[H]
\centering
\caption{Messwerte für die Leckratenmethode der Drehschieberpumpe bei $p_\text{g} = \qty{0,507}{\milli\bar}$.}
\label{tab:11}
\begin{tabular}{rrr}
\toprule
Zeit / s & Druck / mbar & Fehler / mbar \\
\midrule
1 & 0.17 & 0.02 \\
11 & 0.55 & 0.06 \\
21 & 0.61 & 0.06 \\
31 & 0.63 & 0.06 \\
41 & 0.66 & 0.07 \\
51 & 0.68 & 0.07 \\
61 & 0.71 & 0.07 \\
71 & 0.75 & 0.07 \\
81 & 0.78 & 0.08 \\
91 & 0.80 & 0.08 \\
101 & 0.83 & 0.08 \\
111 & 0.86 & 0.09 \\
121 & 0.89 & 0.09 \\
131 & 0.91 & 0.09 \\
141 & 0.94 & 0.09 \\
151 & 0.97 & 0.10 \\
161 & 0.99 & 0.10 \\
171 & 1.00 & 0.10 \\
181 & 1.10 & 0.11 \\
191 & 1.10 & 0.11 \\
\bottomrule
\end{tabular}
\end{table}

\begin{table}[H]
  \centering
\caption{Messwerte für die Leckratenmethode der Drehschieberpumpe bei $p_\text{g} = \qty{10}{\milli\bar}$.}
\label{tab:12} 
\begin{tabular}{rrr}
\toprule
Zeit / s & Druck / mbar & Fehler / mbar \\
\midrule
1 & 6.30 & 0.63 \\
11 & 13.90 & 36.00 \\
21 & 16.70 & 36.00 \\
31 & 19.50 & 36.00 \\
41 & 22.40 & 36.00 \\
51 & 25.50 & 36.00 \\
61 & 28.30 & 36.00 \\
71 & 31.20 & 36.00 \\
81 & 34.00 & 36.00 \\
91 & 37.00 & 36.00 \\
101 & 40.00 & 36.00 \\
111 & 43.10 & 36.00 \\
121 & 46.30 & 36.00 \\
131 & 49.10 & 36.00 \\
141 & 52.30 & 36.00 \\
151 & 55.30 & 36.00 \\
\bottomrule
\end{tabular}
\end{table}

\begin{table}[H]
  \centering
\caption{Messwerte für die Leckratenmethode der Drehschieberpumpe bei $p_\text{g} = \qty{50.07}{\milli\bar}$.}
\label{tab:13} 
\begin{tabular}{rrr}
\toprule
Zeit / s & Druck / mbar & Fehler / mbar \\
\midrule
1 & 56.40 & 36.00 \\
11 & 78.40 & 36.00 \\
21 & 98.60 & 36.00 \\
31 & 118.90 & 36.00 \\
41 & 139.20 & 36.00 \\
51 & 161.10 & 36.00 \\
61 & 181.20 & 36.00 \\
71 & 202.20 & 36.00 \\
81 & 222.50 & 36.00 \\
91 & 242.70 & 36.00 \\
101 & 264.70 & 36.00 \\
111 & 285.00 & 36.00 \\
121 & 306.70 & 36.00 \\
131 & 328.60 & 36.00 \\
141 & 350.40 & 36.00 \\
151 & 372.30 & 36.00 \\
161 & 393.90 & 36.00 \\
171 & 415.40 & 36.00 \\
181 & 437.20 & 36.00 \\
191 & 458.70 & 36.00 \\
\bottomrule
\end{tabular}
\end{table}

\begin{table}[H]
  \caption{Messwerte für die Leckratenmethode der Drehschieberpumpe bei $p_\text{g} = \qty{100}{\milli\bar}$. Zwei Messungen für Mittelwertsbildung.}
\label{tab:14} 
\begin{tabular}{rrrrrr}
\toprule
Zeit / s & Druck 4.1 / mbar & Fehler 4.1 / mbar & Druck 4.2 / mbar & Fehler 4.2 / mbar \\
\midrule
1 & 86.70 & 36.00 & 86.20 & 36.00 \\
11 & 135.10 & 36.00 & 133.50 & 36.00 \\
21 & 174.90 & 36.00 & 172.80 & 36.00 \\
31 & 216.40 & 36.00 & 213.90 & 36.00 \\
41 & 256.10 & 36.00 & 253.10 & 36.00 \\
51 & 295.80 & 36.00 & 292.30 & 36.00 \\
61 & 335.50 & 36.00 & 331.50 & 36.00 \\
71 & 375.30 & 36.00 & 370.70 & 36.00 \\
81 & 417.70 & 36.00 & 412.70 & 36.00 \\
91 & 456.90 & 36.00 & 451.50 & 36.00 \\
101 & 495.60 & 36.00 & 489.80 & 36.00 \\
111 & 536.70 & 36.00 & 530.60 & 36.00 \\
121 & 574.30 & 36.00 & 570.50 & 36.00 \\
131 & 614.00 & 36.00 & 609.60 & 36.00 \\
141 & 652.70 & 36.00 & 647.80 & 36.00 \\
151 & 689.90 & 36.00 & 684.70 & 36.00 \\
161 & 725.50 & 36.00 & 720.00 & 36.00 \\
171 & 759.70 & 36.00 & 753.90 & 36.00 \\
181 & 791.80 & 36.00 & 786.00 & 36.00 \\
191 & 821.90 & 36.00 & 816.10 & 36.00 \\
\bottomrule
\end{tabular}
\end{table}



\subsection{Bestimmung des Saugvermögens der Turbomolekularpumpe durch die Leckratenmethode}
\label{sec:100}
Die Berechnung erfolgt analog zu der bei der Drehschieberpumpe. In \autoref{tab:messung1},
\autoref{tab:messung2},\autoref{tab:messung3},\autoref{tab:messung4} befinden 
sich die Messwerte. Es folgen die geplotteten Messdaten mit den berechneten
Ausgleichsgeraden.

\begin{figure}[H]
  \includegraphics{"TurbomolekularpumpeLeckrate1.pdf"}
  \caption*{}
\end{figure}
\noindent Der Steigungsparameter der Ausgleichsgeraden beträgt bei einem
Gleichgewichtsdruck von $\qty{5.04e-5}{\milli\bar}$
\begin{equation}
  m_{\qty{5.04e-5}{\milli\bar}} = \qty{413(7.3)e-7}{\milli\bar\per\second}.
\end{equation}

\begin{figure}[H]
  \includegraphics{"TurbomolekularpumpeLeckrate2.pdf"}
  \caption*{}
\end{figure}
\noindent Aus der Regression ergibt sich bei $p_g = \qty{56.96e-5}{\milli\bar}$:
\begin{equation}
  m_{\qty{56.96e-5}{\milli\bar}} = \qty{883(20)e-7}{\milli\bar\per\second}.
\end{equation}

\begin{figure}[H]
  \includegraphics{"TurbomolekularpumpeLeckrate3.pdf"}
  \caption*{}
\end{figure}
\noindent Für $p_g = \qty{1.03e-4}{\milli\bar}$ wird ein Parameter von
\begin{equation}
  m_{\qty{1.03e-4}{\milli\bar}} = \qty{883(20)e-7}{\milli\bar\per\second}
\end{equation}
berechnet.

\begin{figure}[H]
  \includegraphics{"TurbomolekularpumpeLeckrate4.pdf"}
  \caption*{}
\end{figure}
\noindent Letztlich ergibt sich für $p_g = \qty{1.91e-4}{\milli\bar}$ eine 
Steigung von
\begin{equation}
  m_{\qty{1.91e-4}{\milli\bar}} = \qty{1580(39)e-7}{\milli\bar\per\second}.
\end{equation}


\noindent Aus diesen Steigungen der Ausgleichsrechnungen wird mittels \autoref{eqn:20}
das Saugvermögen bei den unterschiedlichen drücken bestimmt. Es ergibt sich: 
\begin{align}
  S_\text{eff,\qty{5.04e-5}{\milli\bar}}^\text{leck} =& \qty{27.9+-2.8}{\liter\per\second}\\
  S_\text{eff,\qty{56.96e-5}{\milli\bar}}^\text{leck} =& \qty{43+-4}{\liter\per\second}\\
  S_\text{eff,\qty{1.03e-4}{\milli\bar}}^\text{leck} =& \qty{29.2+-3.0}{\liter\per\second}\\
  S_\text{eff,\qty{1.91e-4}{\milli\bar}}^\text{leck} =& \qty{28.0+-2.9}{\liter\per\second}
\end{align}

\begin{table}[H]
    \centering
    \caption{Messdaten des Druckanstiegs zur Bestimmung der Leckrate bei $p_g = \qty{5.04e-5}{\milli\bar}$.}
    \label{tab:messung1}
    \begin{tabular}{S[table-format=3.0] S[table-format=1.2e-1] @{ \pm } S[table-format=1.2e-1]}
        \toprule
        {Zeit $t / \si{\second}$} & \multicolumn{2}{c}{Druck $p / \si{\milli\bar}$} \\
        \midrule
        1   & 5,00e-5 & 1,50e-5 \\
        11  & 2,06e-4 & 6,18e-5 \\
        21  & 4,63e-4 & 1,39e-4 \\
        31  & 7,31e-4 & 2,19e-4 \\
        41  & 1,12e-3 & 3,36e-4 \\
        51  & 1,50e-3 & 4,50e-4 \\
        61  & 1,90e-3 & 5,70e-4 \\
        71  & 2,50e-3 & 7,50e-4 \\
        81  & 3,08e-3 & 9,24e-4 \\
        91  & 3,64e-3 & 1,09e-3 \\
        101 & 4,38e-3 & 1,31e-3 \\
        111 & 4,95e-3 & 1,49e-3 \\
        121 & 5,36e-3 & 1,61e-3 \\
        131 & 5,99e-3 & 1,80e-3 \\
        \bottomrule
    \end{tabular}
\end{table}

\begin{table}[H]
    \centering
    \caption{Messdaten des Druckanstiegs zur Bestimmung der Leckrate bei $p_g = \qty{56.96e-5}{\milli\bar}$.}
    \label{tab:messung2}
    \begin{tabular}{S[table-format=3.0] S[table-format=1.2e-1] @{ \pm } S[table-format=1.2e-1]}
        \toprule
        {Zeit $t / \si{\second}$} & \multicolumn{2}{c}{Druck $p / \si{\milli\bar}$} \\
        \midrule
        1   & 1,02e-4 & 3,06e-5 \\
        11  & 5,15e-4 & 1,55e-4 \\
        21  & 1,24e-3 & 3,72e-4 \\
        31  & 2,08e-3 & 6,24e-4 \\
        41  & 3,25e-3 & 9,75e-4 \\
        51  & 4,55e-3 & 1,37e-3 \\
        61  & 5,54e-3 & 1,66e-3 \\
        71  & 6,94e-3 & 2,08e-3 \\
        81  & 8,79e-3 & 2,64e-3 \\
        91  & 1,03e-2 & 3,09e-3 \\
        101 & 1,11e-2 & 3,33e-3 \\
        111 & 1,20e-2 & 3,60e-3 \\
        121 & 1,31e-2 & 3,93e-3 \\
        \bottomrule
    \end{tabular}
\end{table}

\begin{table}[H]
    \centering
    \caption{Messdaten des Druckanstiegs zur Bestimmung der Leckrate bei $p_g = \qty{1.03e-4}{\milli\bar}$.}
    \label{tab:messung3}
    \begin{tabular}{S[table-format=3.0] S[table-format=1.2e-1] @{ \pm } S[table-format=1.2e-1]}
        \toprule
        {Zeit $t / \si{\second}$} & \multicolumn{2}{c}{Druck $p / \si{\milli\bar}$} \\
        \midrule
        1   & 1,02e-4 & 3,06e-5 \\
        11  & 5,15e-4 & 1,55e-4 \\
        21  & 1,24e-3 & 3,72e-4 \\
        31  & 2,08e-3 & 6,24e-4 \\
        41  & 3,25e-3 & 9,75e-4 \\
        51  & 4,55e-3 & 1,37e-3 \\
        61  & 5,54e-3 & 1,66e-3 \\
        71  & 6,94e-3 & 2,08e-3 \\
        81  & 8,79e-3 & 2,64e-3 \\
        91  & 1,03e-2 & 3,09e-3 \\
        101 & 1,11e-2 & 3,33e-3 \\
        111 & 1,20e-2 & 3,60e-3 \\
        121 & 1,31e-2 & 3,93e-3 \\
        \bottomrule
    \end{tabular}
\end{table}

\begin{table}[H]
    \centering
    \caption{Messdaten des Druckanstiegs zur Bestimmung der Leckrate bei $p_g = \qty{1.91e-4}{\milli\bar}$.}
    \label{tab:messung4}
    \begin{tabular}{S[table-format=3.0] S[table-format=1.2e-1] @{ \pm } S[table-format=1.2e-1]}
        \toprule
        {Zeit $t / \si{\second}$} & \multicolumn{2}{c}{Druck $p / \si{\milli\bar}$} \\
        \midrule
        1   & 1,90e-4 & 5,70e-5 \\
        11  & 9,61e-4 & 2,88e-4 \\
        21  & 2,43e-3 & 7,29e-4 \\
        31  & 4,48e-3 & 1,34e-3 \\
        41  & 6,32e-3 & 1,90e-3 \\
        51  & 9,25e-3 & 2,78e-3 \\
        61  & 1,11e-2 & 3,33e-3 \\
        71  & 1,27e-2 & 3,81e-3 \\
        81  & 1,46e-2 & 4,38e-3 \\
        91  & 1,68e-2 & 5,04e-3 \\
        101 & 1,90e-2 & 5,70e-3 \\
        111 & 2,10e-2 & 6,30e-3 \\
        121 & 2,28e-2 & 6,84e-3 \\
        131 & 2,47e-2 & 7,41e-3 \\
        141 & 2,67e-2 & 8,01e-3 \\
        151 & 2,89e-2 & 8,67e-3 \\
        \bottomrule
    \end{tabular}
\end{table}

\subsection{Evakuierungskurve der Turbomolekularpumpe}
Analog zur Bestimmung des Saugvermögens der Drehschieberpumpe wird hier 
das Saugvermögen der Turbopumpe durch die Evakuierungskurve bestimmt.
%Die berechnung erfolgt durch \autoref{eqn:100} mit den gemesenen druckwerten $p_\text{E} = \qty{3.52(1)e-6}{\milli\bar}$ und $p_\text{0} = \qty{5e(1.5)-3}{\milli\bar}$, sowie dem 
Das Volumen der Apperatur liegt bei $V = \qty{33+-3.3}{\liter}$.
\noindent Die gemittelten Messwerte, sowie deren Fehler sind \autoref{tab:100}
zu entnehmen. Dabei ergibt sich der Fehler aus dem Fehler des Mittelwerts 
und dem systematischen Fehler. Die geplotteten Messwerte inklusive Ausgleichsgeraden
befinden sich in \autoref{fig:100}. Dabei wurden für die Ausgleichsrechnung die
ersten 12 Messwerte berücksichtigt, da dort ein linearer Abfall erkennbar ist.
Danach nähert sich der Druck dem Enddruck an und es wird kaum noch gepumpt.

\begin{figure}
  \centering
  \includegraphics{build/turboevakuierung.pdf}
  \caption{Darstellung des Logarithmus in Abhängigkeit der Zeit mit Ausgleichsgeraden}
  \label{fig:100}
\end{figure}
\noindent Der Steigungsparameter der Ausgleichsrechnung beträgt 
\begin{equation}
  m = \qty{-46.5(1.4)e-2}{\milli\bar\per\second}.
\end{equation}
Daraus ergibt sich ein Saugvermögen von
\begin{equation}
  S_\text{eff}^\text{Eva} = \qty{15.4+-1.6}{\liter\per\second}.
\end{equation}


\begin{table}[H]
    \centering
    \caption{Messdaten und berechnete Logarithmuswerte}
    \begin{tabular}{c c c}
        \hline
        Zeit $t$ / s & Druck $p$ / mbar & $\ln(\frac{p - p_E}{p_0 - p_E}) \pm \Delta \ln$ \\
        \hline
        0  & $(4,99 \pm 0,50) \cdot 10^{-3}$ & $0,35 \pm 0,10$ \\
        10 & $(5,31 \pm 10,5) \cdot 10^{-5}$ & $-4,27 \pm 2,14$ \\
        20 & $(1,22 \pm 2,44) \cdot 10^{-5}$ & $-6,00 \pm 2,81$ \\
        30 & $(1,01 \pm 2,02) \cdot 10^{-5}$ & $-6,30 \pm 3,09$ \\
        40 & $(7,42 \pm 14,8) \cdot 10^{-6}$ & $-6,91 \pm 4,02$ \\
        50 & $(6,36 \pm 12,7) \cdot 10^{-6}$ & $-7,24 \pm 4,80$ \\
        60 & $(5,30 \pm 10,6) \cdot 10^{-6}$ & $-7,65 \pm 6,10$ \\
        70 & $(5,30 \pm 10,6) \cdot 10^{-6}$ & $-7,65 \pm 6,10$ \\
        80 & $(4,24 \pm 8,48) \cdot 10^{-6}$ & $-8,24 \pm 9,07$ \\
        90 & $(4,24 \pm 8,48) \cdot 10^{-6}$ & $-8,24 \pm 9,07$ \\
        100& $(4,24 \pm 8,48) \cdot 10^{-6}$ & $-8,24 \pm 9,07$ \\
        \hline
    \end{tabular}
    \label{tab:100}
\end{table}

\subsection{Engstelle und Leitwert}
Als letztes wird zwischen den Beiden Pirani- Druckmessgeräten eine 
erheblichere Engstelle eingebaut, sodass sich sofort eine größere Druckdifferenz
zwischen den Sensoren einstellt, welche aus dem geringeren Leitwert der Engstelle
resultiert. Nun werden mit beiden Messgeräten analog zu \autoref{sec:100} die
Saugvermögen berechnet um den Unterschied zu verdeutlichen. Die Messwerte und
Regression für den Drucksensor vor der Engstelle, also zwischen Pumpe und Engstelle 
befinden sich in \autoref{fig:101}, sowie die für das Messgerät nach der Engstelle
(D2) in \autoref{fig:102}. Die Messwerte befinden sich in \autoref{tab:105}.

\begin{figure}[H]
  \caption{Ausgleichsrechnung zur bestimung des Saugvermögens vor der Engstelle}
\label{fig:101}
\includegraphics{build/letzte1.pdf}
\end{figure}

\begin{figure}[H]
  \caption{Ausgleichsrechnung zur bestimung des Saugvermögens hinter der Engstelle}

  \label{fig:102}
\includegraphics{build/letzte2.pdf}
\end{figure}

\noindent Die Parameter der Ausgleichsrechnungen ergeben sich zu 
\begin{align}
m_{\text{vor}}    &= \qty{1.99(0.01)e-04}{\milli\bar\per\second} \\
m_{\text{hinter}} &= \qty{1.84(0.009)e-04}{\milli\bar\per\second}.
\end{align}

\noindent Mithilfe dieser Werte werden die Saugvermögen als 
\begin{align}
S_\text{vor} =& \qty{35.1(3.5)}{\liter\per\second}\\
S_\text{hinter} =& \qty{2.68(1.0)}{\liter\per\second}
\end{align}
bestimmt.

\begin{table}[H]
\centering
\caption{Messdaten und berechnete Logarithmuswerte}
\begin{tabular}{|c|c|c|}
\hline
t / s & p / \unit{\milli\bar} & p/\unit{\milli\bar} \\ \hline
1 & 1.93e-04 $\pm$ 5.79e-05 & 2.34e-03 $\pm$ 7.02e-04 \\ \hline
11 & 2.19e-03 $\pm$ 6.57e-04 & 3.04e-03 $\pm$ 9.12e-04 \\ \hline
21 & 4.22e-03 $\pm$ 1.27e-03 & 5.02e-03 $\pm$ 1.51e-03 \\ \hline
31 & 6.01e-03 $\pm$ 1.80e-03 & 6.88e-03 $\pm$ 2.06e-03 \\ \hline
41 & 9.06e-03 $\pm$ 2.72e-03 & 9.29e-03 $\pm$ 2.79e-03 \\ \hline
51 & 1.09e-02 $\pm$ 3.27e-03 & 1.11e-02 $\pm$ 3.33e-03 \\ \hline
61 & 1.24e-02 $\pm$ 3.72e-03 & 1.26e-02 $\pm$ 3.78e-03 \\ \hline
71 & 1.41e-02 $\pm$ 4.23e-03 & 1.42e-02 $\pm$ 4.26e-03 \\ \hline
81 & 1.61e-02 $\pm$ 4.83e-03 & 1.59e-02 $\pm$ 4.77e-03 \\ \hline
91 & 1.83e-02 $\pm$ 5.49e-03 & 1.78e-02 $\pm$ 5.34e-03 \\ \hline
101 & 2.02e-02 $\pm$ 6.06e-03 & 1.99e-02 $\pm$ 5.97e-03 \\ \hline
111 & 2.18e-02 $\pm$ 6.54e-03 & 2.19e-02 $\pm$ 6.57e-03 \\ \hline
121 & 2.35e-02 $\pm$ 7.05e-03 & 2.33e-02 $\pm$ 6.99e-03 \\ \hline
131 & 2.53e-02 $\pm$ 7.59e-03 & 2.51e-02 $\pm$ 7.53e-03 \\ \hline
141 & 2.72e-02 $\pm$ 8.16e-03 & 2.68e-02 $\pm$ 8.04e-03 \\ \hline
\end{tabular}
\label{tab:105}
\end{table}

\subsection{Vergleich und druckabhängige Darstellung der Saugvermögen beider Pumpen}
Als Zusammenfassung aller berechneter Saugvermögen sollen diese hier in Abhängigkeit
des Drucks zusammengefasst werden. In \autoref{fig:103} befinden sich die Saugvermögen
der Drehschieberpumpe und in \autoref{fig:104} jene der Turbomolekularpumpe.

\begin{figure}[h]
  \caption{Saugvermögen der Drehschieberpumpe in Abhängigkeit des Drucks.}
  \label{fig:103}
\includegraphics{build/saugvermögenDrehschieberpumpe.pdf}
\end{figure}

\begin{figure}[h]
  \caption{Saugvermögen der Turbomolekularpumpe in Abhängigkeit des Drucks.}
  \label{fig:104}
\includegraphics{build/saugvermögenTurbopumpe.pdf}
\end{figure}

\section*{Fehlerrechnung}
% Die Berechnung des Logarithmus erfolgt über die Formel
% \begin{equation}
%     y = \ln\left( \frac{p - p_E}{p_0 - p_E} \right)
% \end{equation}
% \noindent Da sowohl der gemessene Druck $p$ als auch der Enddruck $p_E$
% fehlerbehaftet sind ($\Delta p$ und $\Delta p_E$), berechnet sich der Fehler $\Delta y$ über die Gaußsche Fehlerfortpflanzung
% \begin{equation}
%     \Delta y = \sqrt{ \left( \frac{\partial y}{\partial p} \cdot \Delta p \right)^2 + \left( \frac{\partial y}{\partial p_E} \cdot \Delta p_E \right)^2 }.
% \end{equation}
% Die Ableitungen berechnen sich als
% \begin{equation}
%     \frac{\partial y}{\partial p} = \frac{1}{p - p_E},
% \end{equation}
% \begin{equation}
%     \frac{\partial y}{\partial p_E} = \frac{p - p_0}{(p - p_E)(p_0 - p_E)}.
% \end{equation}
% Einsetzen in die Fehlerfortpflanzung liefert den Fehler des Logarithmus:
% \begin{equation}
%     \Delta y = \sqrt{ \left( \frac{\Delta p}{p - p_E} \right)^2 + \left( \frac{(p - p_0) \cdot \Delta p_E}{(p - p_E)(p_0 - p_E)} \right)^2 }.
% \end{equation}
% Des weiteren gibt es Bei der Bildung von Mittelwerten noch den fehler des Mittelwertes $\delta x$, welcher sich mit dem Systematischen fehler wie folgt zusamensetzt
% \begin{equation}
%   \delta y_\text{gesamt} = \sqrt{(\delta x)^2+(\delta y)^2}
% \end{equation}

Die Berechnung des Logarithmus erfolgt über die Formel
\begin{equation}
    y = \ln\left( \frac{p - p_{\text{E}}}{p_0 - p_{\text{E}}} \right).
\end{equation}
Da sowohl der gemessene Druck $p$ als auch der Enddruck $p_{\text{E}}$ fehlerbehaftet sind ($\Delta p$ und $\Delta p_{\text{E}}$), berechnet sich der Fehler $\Delta y$ über die Gaußsche Fehlerfortpflanzung:
\begin{equation}
    \Delta y = \sqrt{ \left( \frac{\partial y}{\partial p} \cdot \Delta p \right)^2 + \left( \frac{\partial y}{\partial p_{\text{E}}} \cdot \Delta p_{\text{E}} \right)^2 }.
\end{equation}
Die partiellen Ableitungen berechnen sich zu
\begin{align}
    \frac{\partial y}{\partial p}      &= \frac{1}{p - p_{\text{E}}}, \\[6pt]
    \frac{\partial y}{\partial p_{\text{E}}} &= \frac{p - p_0}{(p - p_{\text{E}})(p_0 - p_{\text{E}})}.
\end{align}
Einsetzen in die Fehlerfortpflanzung liefert den Fehler des Logarithmus:
\begin{equation}
    \Delta y = \sqrt{ \left( \frac{\Delta p}{p - p_{\text{E}}} \right)^2 + \left( \frac{(p - p_0) \, \Delta p_{\text{E}}}{(p - p_{\text{E}})(p_0 - p_{\text{E}})} \right)^2 }.
\end{equation}
Bei der Bildung von Mittelwerten kommt zum statistischen Fehler des Mittelwerts
$\delta x$ noch der systematische Fehler $\delta y$ hinzu. Der Gesamtfehler ergibt sich somit zu
\begin{equation}
    \delta y_{\text{ges}} = \sqrt{(\delta x)^2 + (\delta y)^2}.
\end{equation}